\documentclass{article}

\usepackage[french]{babel}
\usepackage[utf8]{inputenc}
\usepackage[T1]{fontenc}
\usepackage{amsfonts}            %For \leadsto
\usepackage{amsmath}             %For \text
\usepackage{fancybox}            %For \ovalbox
\usepackage{color}
\usepackage{geometry}
\usepackage{array}
\usepackage{hyperref}

\DeclareUnicodeCharacter{03BB}{\ensuremath{{\color{black}{\lambda}}}}

\title{Projet \#2}
\author{IFT-3225}

\begin{document}

\maketitle

\newcommand \mML {\ensuremath\mu\textsl{ML}}
\newcommand \kw [1] {\textsf{#1}}
\newcommand \id [1] {\textsl{#1}}
\newcommand \punc [1] {\kw{`#1'}}
\newcommand \str [1] {\texttt{"#1"}}
\newenvironment{outitemize}{
  \begin{itemize}
  \let \origitem \item \def \item {\origitem[]\hspace{-18pt}}
}{
  \end{itemize}
}
\newcommand \Align [2][t] {\begin{array}[#1]{@{}l} #2 \end{array}}

\noindent Yanis Boulogne (20316250)\newline
Karl-Antoine Plouffe (20218864) \newline

\section{Répartition des tâches}
\begin{table}[htbp]
\centering
    \setlength{\leftmargini}{0.5cm}
    \begin{tabular}{| p{5cm} | p{5cm} |}
        \hline
        Yanis Boulogne & Karl-Antoine Plouffe \\
        \hline
        \begin{itemize} 
            \item Chemins d'api back-end
            \item Faire un sequelize maison
            \item Style de la documentation
        \end{itemize} & 
        \begin{itemize} 
            \item Implémentation de la section D- Page Web 
            \item REST relatif à l'application one-page
            \item Style de l'application de jeu
        \end{itemize} \\
        \hline
    \end{tabular}
\end{table}

\section{Page HTML abritée au DIRO}
L'application n'a pas de page HTML dédiée, seuls des modèles EJS existent dans le dossier suivant : \\
\url{https://www-ens.iro.umontreal.ca/hiver/~boulogny/TangoLucioUltimate/src/views/} \\
Ces modèles permettent de d'afficher dynamiquement l'application one-page dans un navigateur web en utilisant les chemins de la section D en localhost lorsque notre application node.js est lancée.

\newpage

\section{Ressources utilisées}
\subsection{Langages}
\begin{itemize}
  \item HTML
  \item JavaScript
  \item PHP
  \item SQL
  \item css
\end{itemize}
\subsection{Librairies}
\begin{itemize}  
  \item Bootstrap 5.3.3
  \item DataTables 2.2.2
  \item Bcrypt
\end{itemize}
\subsection{Frameworks et modules}
\begin{itemize}  
  \item Node.js
  \item EJS 3.1.10
  \item Express 4.21.2
  \item express-session 1.18.1
  \item express-ejs-layouts 2.5.1
  \item Node.js Path
  \item Axios
  \item CORS 
  \item MySQL2
\end{itemize}

\section{Commentaires et précisions}

\subsection{Lancer l'application}
Le projet ayant été fait avec Node.js, il est important avant de le lancer de faire :

\begin{verbatim}
npm install
\end{verbatim}

dans la console afin d'installer les dépendances citées précédemment.

Ensuite, il suffit de faire :

\begin{verbatim}
node server.js &
\end{verbatim}

pour lancer l'application et obtenir son \textit{PID}, utile pour arrêter l'application avec la commande suivante :

\begin{verbatim}
kill {pid}
\end{verbatim}


\subsection{Sequelize et connection à la DB du DIRO}
Nous avions commencé le projet en JavaScript et utilisé le module "sequelize" mais à cause des restrictions du DIRO, il était impossible de se connecter à notre base de données. Au bout d'une semaine, nous nous sommes résolus à utiliser du PHP, mais sans tout refaire de 0, ce qui nous a amenés à créer notre propre module sequelize.
Le fichier db-bridge.php fait exactement cela, les fichiers PHP pouvant communiquer avec les BDD du DIRO, nous avons fait le choix de faire un pont entre toute notre API existante et le DIRO. La manière dont cela fonctionne est plutôt simple, dans chaque modèle que nous avons, nous avons créé des mots-clés qui font appel à database.js pour remettre au propre les appels SQL et le pont PHP renvoie le résultat en format JSON exploitable par le reste de l'API.

\subsection{Joueurs anonymes}
Nous permettons à tous de jouer aux jeux de l'application, même les utilisateurs non authentifiés. Les points gagnés lors de parties ne sont évidemmment pas enregistrés s'il n'y a pas d'utilisateur connecté.

\end{document}